\chapter{Base documental y trazabilidad}
\label{ap:base-documental}

La base funcional de esta memoria es el documento
\emph{Proyecto plataforma SaaS de votación digital: Especificación de
Requisitos}, versión 1.0, de 7 de octubre de 2026. Los identificadores OBJ,
ACT, RF, RI y RNF se conservan para facilitar la trazabilidad con esa fuente.

La especificación completa contiene los flujos normales, precondiciones,
excepciones y criterios de verificación de cada requisito. En esta memoria se
incluyen el catálogo funcional resumido y los elementos necesarios para
entender el análisis, junto con el mapeo de RNF y el cronograma. Las decisiones de diseño y los umbrales propuestos
deben revisarse con el tutor antes de usarlos como criterios finales de
aceptación.

\section{Fuentes PlantUML de los diagramas}
\label{ap:fuentes-plantuml}

Los diagramas de las figuras~\ref{fig:casos-uso-gestion} y
\ref{fig:casos-uso-votacion} se mantienen en los siguientes ficheros
PlantUML, que permiten revisar y regenerar sus imágenes PDF. La arquitectura y la secuencia se mantienen como fuentes TikZ vectoriales; el flujo de secuencia también cuenta con una fuente PlantUML equivalente en el listado A.3.

\lstinputlisting[inputpath=./,caption={Fuente PlantUML de gestión y configuración.},label={lst:casos-uso-gestion}]{Figuras/diagramas/casos_uso_gestion.puml}

\lstinputlisting[inputpath=./,caption={Fuente PlantUML de participación y verificación.},label={lst:casos-uso-votacion}]{Figuras/diagramas/casos_uso_votacion.puml}

\lstinputlisting[inputpath=./,caption={Fuente PlantUML equivalente del flujo de emisión y verificación del recibo.},label={lst:secuencia-voto-plantuml}]{Figuras/diagramas/secuencia_emitir_voto_recibo.puml}

\lstinputlisting[inputpath=./,caption={Fuente PlantUML de actividad para el escrutinio y la verificación.},label={lst:actividad-escrutinio-plantuml}]{Figuras/diagramas/escrutinio_verificacion_actividad.puml}
